% The Roofline model (Lesson 8), log-log axes, for our example machine
% with a peak of 16 TFLOP/s and 1 TB/s of memory bandwidth (ridge point
% 16 FLOP/B).  The slanted roof is bandwidth x intensity (memory-bound
% region), the flat roof the peak (compute-bound region).  Kernel A has
% an arithmetic intensity of 0.5 FLOP/B, kernel B of 2 FLOP/B: raising
% the intensity moves a kernel right, toward the ridge point.
% Redraws the outline figure m8l3_roffline [sic] with our own numbers.
% Usage: % The Roofline model (Lesson 8), log-log axes, for our example machine
% with a peak of 16 TFLOP/s and 1 TB/s of memory bandwidth (ridge point
% 16 FLOP/B).  The slanted roof is bandwidth x intensity (memory-bound
% region), the flat roof the peak (compute-bound region).  Kernel A has
% an arithmetic intensity of 0.5 FLOP/B, kernel B of 2 FLOP/B: raising
% the intensity moves a kernel right, toward the ridge point.
% Redraws the outline figure m8l3_roffline [sic] with our own numbers.
% Usage: % The Roofline model (Lesson 8), log-log axes, for our example machine
% with a peak of 16 TFLOP/s and 1 TB/s of memory bandwidth (ridge point
% 16 FLOP/B).  The slanted roof is bandwidth x intensity (memory-bound
% region), the flat roof the peak (compute-bound region).  Kernel A has
% an arithmetic intensity of 0.5 FLOP/B, kernel B of 2 FLOP/B: raising
% the intensity moves a kernel right, toward the ridge point.
% Redraws the outline figure m8l3_roffline [sic] with our own numbers.
% Usage: % The Roofline model (Lesson 8), log-log axes, for our example machine
% with a peak of 16 TFLOP/s and 1 TB/s of memory bandwidth (ridge point
% 16 FLOP/B).  The slanted roof is bandwidth x intensity (memory-bound
% region), the flat roof the peak (compute-bound region).  Kernel A has
% an arithmetic intensity of 0.5 FLOP/B, kernel B of 2 FLOP/B: raising
% the intensity moves a kernel right, toward the ridge point.
% Redraws the outline figure m8l3_roffline [sic] with our own numbers.
% Usage: \input{figures/l08-roofline}
% Width: body (120 mm); the drawing is about 110 mm wide.
\begin{tikzpicture}[x=8.5mm, y=6mm, font=\footnotesize\sffamily,
    >={Stealth[length=1.5mm]}]
  % regions (x = log2 FLOP/B from -3 to 7, y = log2 TFLOP/s from -3 to 5)
  \fill[ln-accent!8] (-3,-3) rectangle (4,5.4);
  \fill[ln-tip!10] (4,-3) rectangle (7,5.4);
  \node[text=ln-accent, align=center] at (-0.2,4.6)
    {\iflangja{メモリ律速}{memory-bound}};
  \node[text=ln-tip, align=center] at (5.5,0)
    {\iflangja{演算律速}{compute-bound}};
  % axes
  \draw[->] (-3,-3) -- (7.3,-3);
  \draw[->] (-3,-3) -- (-3,5.6);
  \foreach \e/\l in {-3/{1/8},-2/{1/4},-1/{1/2},0/1,1/2,2/4,3/8,4/16,5/32,6/64,7/128} {
    \draw (\e,-3) -- ++(0,-0.15) node[below, inner sep=1pt] {\l};
  }
  \foreach \e/\l in {-3/{1/8},-1/{1/2},1/2,3/8,4/16,5/32} {
    \draw (-3,\e) -- ++(-0.12,0) node[left, inner sep=1pt] {\l};
  }
  \node[below, inner sep=1pt] at (2,-4.1)
    {\iflangja{演算強度（FLOP/B，対数目盛）}{arithmetic intensity (FLOP/B, log scale)}};
  \node[rotate=90, above, inner sep=1pt] at (-4.15,1)
    {\iflangja{性能（TFLOP/s，対数目盛）}{performance (TFLOP/s, log scale)}};
  % roofs
  \draw[very thick, ln-accent] (-3,-3) -- (4,4);
  \draw[very thick, ln-tip] (4,4) -- (7,4);
  \node[above, text=ln-tip] at (5.5,4)
    {\iflangja{ピーク 16\,TFLOP/s}{peak 16\,TFLOP/s}};
  \node[rotate=35.5, below, text=ln-accent] at (2.7,2.7)
    {\iflangja{帯域 1\,TB/s $\times$ 演算強度}{1\,TB/s $\times$ intensity}};
  % ridge point
  \draw[dashed, ln-muted] (4,-3) -- (4,4);
  \fill[ln-caution] (4,4) circle (0.09);
  \node[right, text=ln-caution, align=left] at (4.1,3.25)
    {\iflangja{リッジポイント\\16\,FLOP/B}{ridge point\\16\,FLOP/B}};
  % kernels
  \fill (-1,-1) circle (0.1) node[below right, inner sep=1.5pt] {A: 0.5};
  \fill (1,1) circle (0.1) node[below right, inner sep=1.5pt] {B: 2};
  \draw[->, thick, ln-caution] (-0.85,-0.75) to[bend left=20] (0.8,0.85);
\end{tikzpicture}

% Width: body (120 mm); the drawing is about 110 mm wide.
\begin{tikzpicture}[x=8.5mm, y=6mm, font=\footnotesize\sffamily,
    >={Stealth[length=1.5mm]}]
  % regions (x = log2 FLOP/B from -3 to 7, y = log2 TFLOP/s from -3 to 5)
  \fill[ln-accent!8] (-3,-3) rectangle (4,5.4);
  \fill[ln-tip!10] (4,-3) rectangle (7,5.4);
  \node[text=ln-accent, align=center] at (-0.2,4.6)
    {\iflangja{メモリ律速}{memory-bound}};
  \node[text=ln-tip, align=center] at (5.5,0)
    {\iflangja{演算律速}{compute-bound}};
  % axes
  \draw[->] (-3,-3) -- (7.3,-3);
  \draw[->] (-3,-3) -- (-3,5.6);
  \foreach \e/\l in {-3/{1/8},-2/{1/4},-1/{1/2},0/1,1/2,2/4,3/8,4/16,5/32,6/64,7/128} {
    \draw (\e,-3) -- ++(0,-0.15) node[below, inner sep=1pt] {\l};
  }
  \foreach \e/\l in {-3/{1/8},-1/{1/2},1/2,3/8,4/16,5/32} {
    \draw (-3,\e) -- ++(-0.12,0) node[left, inner sep=1pt] {\l};
  }
  \node[below, inner sep=1pt] at (2,-4.1)
    {\iflangja{演算強度（FLOP/B，対数目盛）}{arithmetic intensity (FLOP/B, log scale)}};
  \node[rotate=90, above, inner sep=1pt] at (-4.15,1)
    {\iflangja{性能（TFLOP/s，対数目盛）}{performance (TFLOP/s, log scale)}};
  % roofs
  \draw[very thick, ln-accent] (-3,-3) -- (4,4);
  \draw[very thick, ln-tip] (4,4) -- (7,4);
  \node[above, text=ln-tip] at (5.5,4)
    {\iflangja{ピーク 16\,TFLOP/s}{peak 16\,TFLOP/s}};
  \node[rotate=35.5, below, text=ln-accent] at (2.7,2.7)
    {\iflangja{帯域 1\,TB/s $\times$ 演算強度}{1\,TB/s $\times$ intensity}};
  % ridge point
  \draw[dashed, ln-muted] (4,-3) -- (4,4);
  \fill[ln-caution] (4,4) circle (0.09);
  \node[right, text=ln-caution, align=left] at (4.1,3.25)
    {\iflangja{リッジポイント\\16\,FLOP/B}{ridge point\\16\,FLOP/B}};
  % kernels
  \fill (-1,-1) circle (0.1) node[below right, inner sep=1.5pt] {A: 0.5};
  \fill (1,1) circle (0.1) node[below right, inner sep=1.5pt] {B: 2};
  \draw[->, thick, ln-caution] (-0.85,-0.75) to[bend left=20] (0.8,0.85);
\end{tikzpicture}

% Width: body (120 mm); the drawing is about 110 mm wide.
\begin{tikzpicture}[x=8.5mm, y=6mm, font=\footnotesize\sffamily,
    >={Stealth[length=1.5mm]}]
  % regions (x = log2 FLOP/B from -3 to 7, y = log2 TFLOP/s from -3 to 5)
  \fill[ln-accent!8] (-3,-3) rectangle (4,5.4);
  \fill[ln-tip!10] (4,-3) rectangle (7,5.4);
  \node[text=ln-accent, align=center] at (-0.2,4.6)
    {\iflangja{メモリ律速}{memory-bound}};
  \node[text=ln-tip, align=center] at (5.5,0)
    {\iflangja{演算律速}{compute-bound}};
  % axes
  \draw[->] (-3,-3) -- (7.3,-3);
  \draw[->] (-3,-3) -- (-3,5.6);
  \foreach \e/\l in {-3/{1/8},-2/{1/4},-1/{1/2},0/1,1/2,2/4,3/8,4/16,5/32,6/64,7/128} {
    \draw (\e,-3) -- ++(0,-0.15) node[below, inner sep=1pt] {\l};
  }
  \foreach \e/\l in {-3/{1/8},-1/{1/2},1/2,3/8,4/16,5/32} {
    \draw (-3,\e) -- ++(-0.12,0) node[left, inner sep=1pt] {\l};
  }
  \node[below, inner sep=1pt] at (2,-4.1)
    {\iflangja{演算強度（FLOP/B，対数目盛）}{arithmetic intensity (FLOP/B, log scale)}};
  \node[rotate=90, above, inner sep=1pt] at (-4.15,1)
    {\iflangja{性能（TFLOP/s，対数目盛）}{performance (TFLOP/s, log scale)}};
  % roofs
  \draw[very thick, ln-accent] (-3,-3) -- (4,4);
  \draw[very thick, ln-tip] (4,4) -- (7,4);
  \node[above, text=ln-tip] at (5.5,4)
    {\iflangja{ピーク 16\,TFLOP/s}{peak 16\,TFLOP/s}};
  \node[rotate=35.5, below, text=ln-accent] at (2.7,2.7)
    {\iflangja{帯域 1\,TB/s $\times$ 演算強度}{1\,TB/s $\times$ intensity}};
  % ridge point
  \draw[dashed, ln-muted] (4,-3) -- (4,4);
  \fill[ln-caution] (4,4) circle (0.09);
  \node[right, text=ln-caution, align=left] at (4.1,3.25)
    {\iflangja{リッジポイント\\16\,FLOP/B}{ridge point\\16\,FLOP/B}};
  % kernels
  \fill (-1,-1) circle (0.1) node[below right, inner sep=1.5pt] {A: 0.5};
  \fill (1,1) circle (0.1) node[below right, inner sep=1.5pt] {B: 2};
  \draw[->, thick, ln-caution] (-0.85,-0.75) to[bend left=20] (0.8,0.85);
\end{tikzpicture}

% Width: body (120 mm); the drawing is about 110 mm wide.
\begin{tikzpicture}[x=8.5mm, y=6mm, font=\footnotesize\sffamily,
    >={Stealth[length=1.5mm]}]
  % regions (x = log2 FLOP/B from -3 to 7, y = log2 TFLOP/s from -3 to 5)
  \fill[ln-accent!8] (-3,-3) rectangle (4,5.4);
  \fill[ln-tip!10] (4,-3) rectangle (7,5.4);
  \node[text=ln-accent, align=center] at (-0.2,4.6)
    {\iflangja{メモリ律速}{memory-bound}};
  \node[text=ln-tip, align=center] at (5.5,0)
    {\iflangja{演算律速}{compute-bound}};
  % axes
  \draw[->] (-3,-3) -- (7.3,-3);
  \draw[->] (-3,-3) -- (-3,5.6);
  \foreach \e/\l in {-3/{1/8},-2/{1/4},-1/{1/2},0/1,1/2,2/4,3/8,4/16,5/32,6/64,7/128} {
    \draw (\e,-3) -- ++(0,-0.15) node[below, inner sep=1pt] {\l};
  }
  \foreach \e/\l in {-3/{1/8},-1/{1/2},1/2,3/8,4/16,5/32} {
    \draw (-3,\e) -- ++(-0.12,0) node[left, inner sep=1pt] {\l};
  }
  \node[below, inner sep=1pt] at (2,-4.1)
    {\iflangja{演算強度（FLOP/B，対数目盛）}{arithmetic intensity (FLOP/B, log scale)}};
  \node[rotate=90, above, inner sep=1pt] at (-4.15,1)
    {\iflangja{性能（TFLOP/s，対数目盛）}{performance (TFLOP/s, log scale)}};
  % roofs
  \draw[very thick, ln-accent] (-3,-3) -- (4,4);
  \draw[very thick, ln-tip] (4,4) -- (7,4);
  \node[above, text=ln-tip] at (5.5,4)
    {\iflangja{ピーク 16\,TFLOP/s}{peak 16\,TFLOP/s}};
  \node[rotate=35.5, below, text=ln-accent] at (2.7,2.7)
    {\iflangja{帯域 1\,TB/s $\times$ 演算強度}{1\,TB/s $\times$ intensity}};
  % ridge point
  \draw[dashed, ln-muted] (4,-3) -- (4,4);
  \fill[ln-caution] (4,4) circle (0.09);
  \node[right, text=ln-caution, align=left] at (4.1,3.25)
    {\iflangja{リッジポイント\\16\,FLOP/B}{ridge point\\16\,FLOP/B}};
  % kernels
  \fill (-1,-1) circle (0.1) node[below right, inner sep=1.5pt] {A: 0.5};
  \fill (1,1) circle (0.1) node[below right, inner sep=1.5pt] {B: 2};
  \draw[->, thick, ln-caution] (-0.85,-0.75) to[bend left=20] (0.8,0.85);
\end{tikzpicture}
